% GENERATED FILE. Edit canonical lessons, then run npm run generate:books.
% canonical-source-hash: fnv1a64:fa6206bd2614712b
% canonical-lessons: HI-C297-phisalna, HI-C297-thokar, HI-C297-larkharana, HI-C297-chilna, HI-C297-lahuluhan

\chapter{A Slip and a Fall}
\label{ch:hi-a-slip-and-a-fall}
\hlchaptermodality{297}

\begin{chapteropening}
\textbf{By the end of this chapter:} \emph{I can say how someone slipped and got hurt.}

Complete the last lesson of chapter 297: I can say how someone slipped and got hurt.
\end{chapteropening}

\section[\texorpdfstring{\hi{फिसलना} (phisalnā)}{फिसलना (phisalnā)}]{\hi{फिसलना} (phisalnā) — to slip}

\label{lesson:HI-C297-phisalna}



\begin{quote}
Before the new one: say the Hindi for a wheel, then the Hindi for to overturn.
\end{quote}

\subsection*{You'll want to know: \hi{फिसलना}}
\textbf{\hi{फिसलना}} — \emph{phisalnā} — \textquotedblleft{}to slip\textquotedblright{}.

\begin{quote}
\hi{बारिश} \hi{के} \hi{बाद} \hi{मीरा} \hi{सीढ़ी} \hi{पर} \hi{फिसल} \hi{गई।}
\end{quote}

\begin{quote}
\emph{bāriś ke bād Mīrā sīṛhī par phisal gaī.}
\end{quote}

\begin{quote}
After the rain Mira slipped on the stairs.
\end{quote}

Past (she …): \hi{वह} \hi{फिसल} \hi{गई}, \emph{vah phisal gaī}.

\begin{cousinweb}
\textbf{\hi{फिसलना}} (\emph{phisalnā}) is inherited from \emph{*phiss}, a form CDIAL reconstructs for slipping (CDIAL 9080). A slippery floor is one you \textbf{\hi{फिसलते} \hi{हैं}} on.
\end{cousinweb}

\begin{grammarlens}[title={What you've built}]
One more word for telling what happened.
\end{grammarlens}

\subsection*{Your turn}
\begin{itemize}
\raggedright
  \item \emph{Say it:} \emph{phisalnā}
  \item \emph{Say it:} \emph{phisalnā}, once more
  \item \emph{Say it:} \emph{bāriś ke bād Mīrā sīṛhī par phisal gaī.} — After the rain Mira slipped on the stairs.
  \item \emph{Recall:} say the Hindi for a wheel, then the Hindi for to overturn, then say \emph{phisalnā} again
\end{itemize}

\begin{tcolorbox}[breakable,colback=teal!4,colframe=teal!35!black,title={Before you move on}]
What does \hi{फिसलना} mean? (\textquotedblleft{}To slip\textquotedblright{}.) Say it once more.
\end{tcolorbox}
\section[\texorpdfstring{\hi{ठोकर} (ṭhokar)}{ठोकर (ṭhokar)}]{\hi{ठोकर} (ṭhokar) — a stumble; a knock}

\label{lesson:HI-C297-thokar}



\begin{quote}
Before the new one: say the Hindi for to overturn, then the Hindi for to slip.
\end{quote}

\subsection*{You'll want to know: \hi{ठोकर}}
\textbf{\hi{ठोकर}} — \emph{ṭhokar} — \textquotedblleft{}a stumble; a knock\textquotedblright{}.

\begin{quote}
\hi{अरुण} \hi{को} \hi{पत्थर} \hi{से} \hi{ठोकर} \hi{लगी।}
\end{quote}

\begin{quote}
\emph{Aruṇ ko patthar se ṭhokar lagī.}
\end{quote}

\begin{quote}
Arun tripped over a stone.
\end{quote}

Feminine.

\begin{cousinweb}
\textbf{\hi{ठोकर}} (\emph{ṭhokar}) is inherited from \emph{*ṭhōkk}, a knock, a form CDIAL reconstructs (CDIAL 5513). The person who trips takes \textbf{\hi{को}}: \textbf{\hi{अरुण} \hi{को} \hi{ठोकर} \hi{लगी}}, a knock struck Arun.
\end{cousinweb}

\begin{grammarlens}[title={What you've built}]
One more word for telling what happened.
\end{grammarlens}

\subsection*{Your turn}
\begin{itemize}
\raggedright
  \item \emph{Say it:} \emph{ṭhokar}
  \item \emph{Say it:} \emph{ṭhokar}, once more
  \item \emph{Say it:} \emph{Aruṇ ko patthar se ṭhokar lagī.} — Arun tripped over a stone.
  \item \emph{Recall:} say the Hindi for to overturn, then the Hindi for to slip, then say \emph{ṭhokar} again
\end{itemize}

\begin{tcolorbox}[breakable,colback=teal!4,colframe=teal!35!black,title={Before you move on}]
What does \hi{ठोकर} mean? (\textquotedblleft{}A stumble; a knock\textquotedblright{}.) Say it once more.
\end{tcolorbox}
\section[\texorpdfstring{\hi{लड़खड़ाना} (laṛkhaṛānā)}{लड़खड़ाना (laṛkhaṛānā)}]{\hi{लड़खड़ाना} (laṛkhaṛānā) — to stagger}

\label{lesson:HI-C297-larkharana}



\begin{quote}
Before the new one: say the Hindi for to slip, then the Hindi for a stumble.
\end{quote}

\subsection*{You'll want to know: \hi{लड़खड़ाना}}
\textbf{\hi{लड़खड़ाना}} — \emph{laṛkhaṛānā} — \textquotedblleft{}to stagger\textquotedblright{}.

\begin{quote}
\hi{ठोकर} \hi{के} \hi{बाद} \hi{वह} \hi{लड़खड़ाया} \hi{और} \hi{गिर} \hi{गया।}
\end{quote}

\begin{quote}
\emph{ṭhokar ke bād vah laṛkhaṛāyā aur gir gayā.}
\end{quote}

\begin{quote}
After the stumble he staggered and fell.
\end{quote}

Past (he …): \hi{वह} \hi{लड़खड़ाया}, \emph{vah laṛkhaṛāyā}.

\begin{cousinweb}
\textbf{\hi{लड़खड़ाना}} (\emph{laṛkhaṛānā}) is set by CDIAL under Sanskrit \emph{laṭati}, \textquotedblleft{}shakes\textquotedblright{} (CDIAL 10916), where Hindi appears as \emph{larkharānā}, to stumble. The repeated \emph{-khaṛ-} sound gives the unsteady steps.
\end{cousinweb}

\begin{grammarlens}[title={What you've built}]
One more word for telling what happened.
\end{grammarlens}

\subsection*{Your turn}
\begin{itemize}
\raggedright
  \item \emph{Say it:} \emph{laṛkhaṛānā}
  \item \emph{Say it:} \emph{laṛkhaṛānā}, once more
  \item \emph{Say it:} \emph{ṭhokar ke bād vah laṛkhaṛāyā aur gir gayā.} — After the stumble he staggered and fell.
  \item \emph{Recall:} say the Hindi for to slip, then the Hindi for a stumble, then say \emph{laṛkhaṛānā} again
\end{itemize}

\begin{tcolorbox}[breakable,colback=teal!4,colframe=teal!35!black,title={Before you move on}]
What does \hi{लड़खड़ाना} mean? (\textquotedblleft{}To stagger\textquotedblright{}.) Say it once more.
\end{tcolorbox}
\section[\texorpdfstring{\hi{छिलना} (chilnā)}{छिलना (chilnā)}]{\hi{छिलना} (chilnā) — to be grazed; to be scraped}

\label{lesson:HI-C297-chilna}



\begin{quote}
Before the new one: say the Hindi for a stumble, then the Hindi for to stagger.
\end{quote}

\subsection*{You'll want to know: \hi{छिलना}}
\textbf{\hi{छिलना}} — \emph{chilnā} — \textquotedblleft{}to be grazed; to be scraped\textquotedblright{}.

\begin{quote}
\hi{गिरने} \hi{से} \hi{उसका} \hi{घुटना} \hi{छिल} \hi{गया।}
\end{quote}

\begin{quote}
\emph{girne se uskā ghuṭnā chil gayā.}
\end{quote}

\begin{quote}
He grazed his knee when he fell.
\end{quote}

Past (it …): \hi{वह} \hi{घुटना} \hi{छिल} \hi{गया}, \emph{vah ghuṭnā chil gayā}.

\begin{cousinweb}
\textbf{\hi{छिलना}} (\emph{chilnā}) is the intransitive of \emph{chīlnā}, to peel, which CDIAL 5051 lists for Hindi under \emph{*chilla}: when you fall, the skin \textbf{\hi{छिल} \hi{जाती} \hi{है}}, peels off by itself.
\end{cousinweb}

\begin{grammarlens}[title={What you've built}]
One more word for telling what happened.
\end{grammarlens}

\subsection*{Your turn}
\begin{itemize}
\raggedright
  \item \emph{Say it:} \emph{chilnā}
  \item \emph{Say it:} \emph{chilnā}, once more
  \item \emph{Say it:} \emph{girne se uskā ghuṭnā chil gayā.} — He grazed his knee when he fell.
  \item \emph{Recall:} say the Hindi for a stumble, then the Hindi for to stagger, then say \emph{chilnā} again
\end{itemize}

\begin{tcolorbox}[breakable,colback=teal!4,colframe=teal!35!black,title={Before you move on}]
What does \hi{छिलना} mean? (\textquotedblleft{}To be grazed; to be scraped\textquotedblright{}.) Say it once more.
\end{tcolorbox}
\section[\texorpdfstring{\hi{लहूलुहान} (lahūluhān)}{लहूलुहान (lahūluhān)}]{\hi{लहूलुहान} (lahūluhān) — covered in blood}

\label{lesson:HI-C297-lahuluhan}



\begin{quote}
Before the new one: say the Hindi for to stagger, then the Hindi for to be grazed.
\end{quote}

\subsection*{You'll want to know: \hi{लहूलुहान}}
\textbf{\hi{लहूलुहान}} — \emph{lahūluhān} — \textquotedblleft{}covered in blood\textquotedblright{}.

\begin{quote}
\hi{गिरने} \hi{के} \hi{बाद} \hi{उसका} \hi{हाथ} \hi{लहूलुहान} \hi{था।}
\end{quote}

\begin{quote}
\emph{girne ke bād uskā hāth lahūluhān thā.}
\end{quote}

\begin{quote}
After the fall his hand was covered in blood.
\end{quote}

\begin{cousinweb}
\textbf{\hi{लहूलुहान}} (\emph{lahūluhān}) is \emph{lahū}, blood, which CDIAL 11165 lists for Hindi under Sanskrit \emph{lohita}, \textquotedblleft{}red\textquotedblright{}, followed by an expressive echo, \emph{-luhān}, that makes the picture stronger. It describes a person or a body part: \textbf{\hi{हाथ} \hi{लहूलुहान} \hi{था}}.
\end{cousinweb}

\begin{grammarlens}[title={What you've built}]
One more word for telling what happened.
\end{grammarlens}

\subsection*{Your turn}
\begin{itemize}
\raggedright
  \item \emph{Say it:} \emph{lahūluhān}
  \item \emph{Say it:} \emph{lahūluhān}, once more
  \item \emph{Say it:} \emph{girne ke bād uskā hāth lahūluhān thā.} — After the fall his hand was covered in blood.
  \item \emph{Recall:} say the Hindi for to stagger, then the Hindi for to be grazed, then say \emph{lahūluhān} again
\end{itemize}

\begin{tcolorbox}[breakable,colback=teal!4,colframe=teal!35!black,title={Before you move on}]
What does \hi{लहूलुहान} mean? (\textquotedblleft{}Covered in blood\textquotedblright{}.) Say it once more.
\end{tcolorbox}
